%------------------------------------------------------------------------------%
% \chapter{Plano Cartesiano e Equações}
%------------------------------------------------------------------------------%

%------------------------------------------------------------------------------%
\begin{exercicios}{Equação da Reta}
    
    \begin{exercicio}
        Encontre a equação da reta que passa pelo ponto $(1,3)$ com declividade~2. 
        Esboce seu gráfico e verifique se os pontos $(2,6)$, $(0,0)$ e $(-1,1)$ pertencem a reta.
    \end{exercicio}

    \begin{exercicio}
        Esboce o gráfico da equação $3x - 4y = 24$.
    \end{exercicio}

    \begin{exercicio}
        Dado que o ponto $(2,-3)$ pertence a reta $-2x+ky+10=0$, determine o valor de $k$.
    \end{exercicio}
    
    \begin{exercicio}
        Encontre as inclinações das retas mostradas na figura abaixo.
        \addgraphic{.6}{exercicio-retas-01}
    \end{exercicio}
        
    \begin{exercicio}
        Determine as inclinações das retas que passam pelos pares de ponto abaixo.
        \begin{alphalist}[2]
            \item $(4,1)$ e $(2,3)$
            \item $(1,2)$ e $(-2,-4)$
            \item $(-5,-2)$ e $(3,-2)$
            \item $(-4,5)$ e $(1,-10)$
            \item $(6,4)$ e $(-3,1)$
            \item $(-3,4)$ e $(7,2)$
            \item $(-2,-6)$ e $(-1,1)$
            \item $(5,-2)$ e $(-9,5)$
        \end{alphalist}
        \begin{answers}
            \begin{mathlist}[4]
                \item m = -1
                \item m = 2
                \item m = 0
                \item m = -3
                \item m = \frac{1}{3}
                \item m = - \frac{1}{5}
                \item m = 7
                \item m = - \frac{1}{2}
            \end{mathlist}
        \end{answers}
    \end{exercicio}
    
    \begin{exercicio}
        Escreva as equações das retas definidas pelas inclinações e interceptos abaixo.
        \begin{alphalist}
            \item Intercepto-$y$:  -1;          inclinação: \sfrac{4}{5}
            \item Intercepto-$y$:   2;          inclinação: \sfrac{-3}{4}
            \item Intercepto-$y$:   4;          inclinação: -3
            \item Intercepto-$y$:  -3;          inclinação: \sfrac{1}{3}
            \item Intercepto-$y$: \sfrac{1}{2}; inclinação:  2
            \item Intercepto-$y$:   0;          inclinação: -1
        \end{alphalist}
        \begin{answers}
            \begin{mathlist}[2]
                \item y = \sfrac{ 4x}{5}    -1
                \item y = \sfrac{-3x}{4}    +2
                \item y = -3x      +4
                \item y = \sfrac{ x}{3}    -3
                \item y =  2x    +\sfrac{1}{2}
                \item y = -x                        
            \end{mathlist}
        \end{answers}
    \end{exercicio}
    
    \begin{exercicio}
        Determine as equações das retas que satisfazem as condições indicadas. Em seguida, trace seus gráficos.
        \begin{alphalist}
            \item Passa por (2,-1)  e tem inclinação 3
            \item Passa por (1,5)   e tem inclinação -3
            \item Passa por (-2,1)  e tem inclinação \sfrac{1}{3}
            \item Passa por (6,-4)  e tem inclinação \sfrac{-1}{3}
            \item Passa por (-4,8)  e tem inclinação -2
            \item Passa por (-3,-2) e tem inclinação \sfrac{3}{2}
        \end{alphalist}
        \begin{answers}
            \begin{mathlist}[3]
                \item y = 3 x - 7
                \item y = - 3 x + 8
                \item y = \sfrac{x}{3} + \sfrac{5}{3}
                \item y = - \sfrac{x}{3} - 2
                \item y = - 2 x
                \item y = \sfrac{3 x}{2} + \sfrac{5}{2}
            \end{mathlist}
        \end{answers}
    \end{exercicio}

    \begin{exercicio}
        Escreva a equação reduzida da reta $r$ que passa pelos pontos $A(-1,1)$ e $B(7,25)$.
    \end{exercicio}
    
    \begin{exercicio}
        Escreva a equação geral das retas que passam pelos pontos
        \begin{alphalist}
            \item $(-2,0)$ \tabto{3cm} $(3,0)$
            \item $(0,-1)$ \tabto{3cm} $(2,0)$
            \item $(-\sfrac{4}{3},0)$ \tabto{3cm} $(0,-2)$
        \end{alphalist}
        \begin{answers}
            \begin{mathlist}[3]
                \item  y + 2  = 0
                \item 2y -  x = 0
                \item 6y + 9x + 20 = 0
            \end{mathlist}
        \end{answers}
    \end{exercicio}
    
    \begin{exercicio}
        Encontre as equações das retas que satisfazem as condições indicadas.
        \begin{alphalist}
            \item Passa por (-1,-3) e intercepta o eixo-$y$ na ordenada 1
            \item Passa por (1,2) e por (2,1)
            \item Passa por (4,-2) e por (-3,-2)
            \item Intercepta o eixo-$y$ na ordenada 3 e o eixo-$x$ na abscissa -2
            \item Intercepta o eixo-$y$ na ordenada 2 e o eixo-$x$ na abscissa 1
            \item Passa por (-2,-6) e intercepta o eixo-$x$ na abscissa 10
            \item Passa por (-6,4) e por (2,-1)
            \item Passa por (-2,8) e por (3,4)
            \item Passa por (-5,0) e por (0,-5)
            \item Passa por (3,14) e por (-4,-28)
            \item Passa por (-1,2 1 ) e por (2,19)
            \item Passa por (3,5;-0,6) e por (-2;2,7)
        \end{alphalist}
    \end{exercicio}
        
    \begin{exercicio}
        Encontre a reta definida pelos pontos 
        \( A\left( \sfrac{7}{2}, \sfrac{5}{2}\right) \) e  
        \( B\left(-\sfrac{5}{2},-\sfrac{7}{7}\right) \).
    \end{exercicio}
        
    \begin{exercicio}
        A reta determinada por $A(a,0)$ e $B(0,b)$ passa por $C(3,4)$. Qual a relação entre $a$ e $b$?
    \end{exercicio}
        
    \begin{exercicio}
        A reta determinada por $A(q,q)$ e $B(3,-2)$ passa pela origem. Qual a relação entre $p$ e $q$?
    \end{exercicio}
    
    \begin{exercicio}
        Desenhe no plano cartesiano as retas cujas equações são
        \begin{mathlist}[2]
            \item y=2x
            \item x+y+3=0
            \item x+y=5
            \item 2y+x=0
            \item x-y+5=0
            \item x-y-4=0
        \end{mathlist}
    \end{exercicio}
    
    \begin{exercicio}
        Determine as equações das retas mostradas na figura abaixo.
        \addgraphic{.6}{exercicio-retas-02}
    \end{exercicio}
    
    \begin{exercicio}
        Trace os gráficos das equações abaixo.
	    \begin{mathlist}[2]
	    	\item y = - \frac{2}{3}x + 1
	    	\item y = 5x - 2
	    	\item y = -2x
	    	\item y = 4
	    	\item x - y = -3
	    	\item 3y - x + 4 = 0
	    \end{mathlist}        
    \end{exercicio}
    
    \begin{exercicio}
        Encontre as equações das retas que satisfazem as condições abaixo. Em seguida, trace os gráficos correspondentes.
        \begin{alphalist}
            \item Passa por (-2,1) e (-2, 5).
            \item Passa por (-7,8) e (10, 8).
            \item Reta vertical que passa por (3,-1).
            \item Reta horizontal que passa por (6,-4).
        \end{alphalist}
    \end{exercicio}
    
    \begin{exercicio}
        Dados os pontos $A(-1,2)$ e $B(3,-1)$
        \begin{alphalist}
            \item Marque os pontos no plano Cartesiano, considerando as abscissas no intervalo $[-3,5]$ e as ordenadas em $[-2, 3]$.
            \item Determine a equação da reta que passa pelos pontos. Trace essa reta no gráfico.
            \item Determine a ordenada do ponto dessa reta no qual a abscissa vale 1.
            \item Determine a abscissa do ponto da reta que tem ordenada 0.
        \end{alphalist}
    \end{exercicio}
    
\end{exercicios}

%------------------------------------------------------------------------------%
